\section{Adversarial Controls and the Transfer Boundary}
\label{sec:submission-checks}

\subsection{A Simpler Physical Competitor}
A stronger simplicity check retains one committed physical and instrumental
state, rather than resetting to equilibrium after each event. Short GLR
alternatives test the same admitted load and line operators over at most
19 reports. Load ON severity is profiled within the original bounds; an
accepted OFF changes the operator while retaining the inferred state.
Alternative histories then expire. This is a complete simpler procedure,
not a deliberately memoryless reset baseline.

The single-state procedure and nominal history bank share fitted nominal
operators; the five-model bank additionally mixes the declared resistance
and reactance alternatives. All use a six-report source-stability clock.
The two archived history-bank
outputs are republished with that clock. The GLR includes a fixed penalty
for its maximum search count; neither this penalty nor a numerically
shared threshold makes its scores posterior probabilities or equalizes
false-alarm rates. Table~\ref{tab:single-state-control} aggregates all twelve records.
% Generated from sealed predictions and successor metrics.
\begin{table}[t]\centering
\caption{Open-development comparison with the same six-report publication clock. All 54 event replays are retained. Threshold score semantics differ; this does not establish equal false-alarm operating points. Missing closure includes missed events.}
\label{tab:single-state-control}
\small\setlength{\tabcolsep}{3pt}
\begin{tabular}{lrrrr}\hline
Procedure & Timely & Extras & Early OFF & No OFF\\\hline
Single state, short GLR & 37/54 & 37 & 0 & 4 \\
Nominal history bank & 44/54 & 6 & 0 & 10 \\
Five-model history bank & 51/54 & 3 & 0 & 3 \\
\hline\end{tabular}
\end{table}

The five-model bank obtains \SSMixTimely{}/\SSMixTotal{} timely decisions
with \SSMixExtras{} extras, versus \SSSimpleTimely{}/\SSSimpleTotal{} and
\SSSimpleExtras{} for the single-state procedure. This favors the full
procedure in this development contrast, but does not isolate history
averaging from severity integration, search or regime selection. These are
three event calendars with repeated conditions. The zero alarms over
\SSNormalSeconds{} normal seconds are too little exposure to establish a
low operational alarm rate. The changed counts relative to the original
publication policy also show that the decision clock is consequential.
The corresponding orphan-restoration counts are \SSSimpleOrphan{},
\SSNomOrphan{} and \SSMixOrphan{} for the single-state, nominal and five-model
procedures; zero premature matched closures does not imply error-free restoration.
The original publisher gave \ContMixTimely{}/\ContMixTotal{} with
\ContMixExtras{} extras for the mixture and
\ContNomTimely{}/\ContNomTotal{} with \ContNomExtras{} for nominal.
The common publisher gives \SSMixTimely{}/\SSMixTotal{} with
\SSMixExtras{} and \SSNomTimely{}/\SSNomTotal{} with \SSNomExtras{},
respectively. These are different decision policies on the same H04 outputs,
not conflicting estimates of one endpoint.

\subsection{What Grouped SIM Training Establishes}
We also fit two ExtraTrees variants using dependency groups, rather than
randomly splitting correlated reports. The approximately 70/30 split is
subdivided into \MLTrainGroups{} fitting, \MLValGroups{} validation and
\MLHoldGroups{} reserved groups. SHORT is selected on validation and
frozen before the reserved evaluation. Table~\ref{tab:grouped-learning}
retains its marked load/line asymmetry. Its
\MLHoldTimely{}/\MLHoldEvents{} timely reserved decisions and
\MLHoldExtras{} extras do not establish adequate general transfer.
The corpus was already open and all physical sources remain represented
in development; this is neither a blind test nor an unseen-source split.
Only one reserved normal dependency group underlies the
\MLHoldNormalSeconds{} normal seconds with \MLHoldNormalAlarms{} alarms.
% Generated by import_s01_learning.py; unequal cohorts are not a paired contrast.
\begin{table}[t]\centering
\caption{Grouped SIM development. SHORT is selected only on validation; the reserved 30\% is evaluated once. These cohorts are distinct from H04.}
\label{tab:grouped-learning}
\small\setlength{\tabcolsep}{4pt}
\begin{tabular}{lrrr}\hline
Variant and partition & Load & Line & Extras\\\hline
Baseline, validation & 5/96 & 87/96 & 0 \\
SHORT, validation & 16/96 & 92/96 & 0 \\
SHORT, reserved & 28/288 & 274/288 & 3 \\
\hline\end{tabular}
\end{table}


Replaying the selected classifier on the common H04 cohort gives
\ETCommonTimely{}/\ETCommonTotal{} timely diagnoses and \ETCommonExtras{}
extras: \ETCommonLoad{}/\ETCommonLoadTotal{} loads and
\ETCommonLine{}/\ETCommonLineTotal{} lines. The physical bank's benefit
therefore lies chiefly in loads in this contrast; the classifier resolves
all tested lines. Its frozen publisher uses six identical argmax labels,
without the physical procedures' score threshold. This is a common-cohort,
common-metric comparison of complete procedures, with different publishers.
H04's normal parent also appears in the reserved SIM fraction and is not
additional independent normal evidence.\AIDraftingCredit{}

\subsection{Measurement Integrity and RAW0001}
The frozen physical filter does not admit RAW0001: the initial prefix
violates its ROCOF noise-scale and temporal-stationarity support. The
following ZIP and segmented-ML replays therefore examine complementary
integrity and source interfaces; they do not validate R11 transfer.
The submitted multi-family artifact is replayed on all
\SSRawWindows{} fixed windows of RAW0001. Its predictions are emitted at
window completion. A separate availability monitor reports missing PMUs
from received flags and finite-value masks; it keeps missing-PMU identity,
corrupt-PMU identity and physical source distinct. Such diagnoses do not
establish malicious intent. Complete-window corruption diagnoses describe
their observed interval and are not broadcast backwards as causal labels.

Against the current annotation support, the submitted artifact detects an
abnormality in \SSRawDetected{}/\SSRawPositive{} positive windows, with
\SSRawFalse{} false-positive window among \SSRawNormal{} normal windows;
all \SSRawAvailable{} windows produce outputs. Correct physical-family
coverage is \SSRawPhysicalCorrect{}/\SSRawPhysicalTotal{}. The availability
monitor identifies all \SSMissingTotal{} missing reports and their PMU,
whereas corruption localization is correct in
\SSCorruptLocated{}/\SSCorruptTotal{} positive windows. A separate spurious
corruption window is retained. These denominators describe one retrospective
record, not independent field trials. Both implementations identify all
\MissingEpisodes{} missing intervals and their PMU; the immediate interface
reduces mean first-notice delay from \MissingLegacyDelay{} to
\MissingImmediateDelay{} seconds on the received-report clock. This limited
publication benefit does not measure wall-clock latency or physical inference.
Its event-code annotations do not certify individual physical actuation
boundaries, so actuation delay and restoration are not scored.

The archived source spreadsheet permits a separate metadata-based source
comparison. A frozen hybrid uses the SIM-trained SHORT classifier only to
replace the submitted artifact's load/line source when a first same-family
diagnosis has already been emitted within that window. Predicted mixed
physical/missing cases are eligible under the same rule. All classes,
integrity diagnoses and window emission times remain fixed. Missing input
splits the ML stream into complete segments with explicit unavailable
startup; their prefixes are not verified normal. This is a segmented ML
transfer test, not a successful execution of the physical history bank.
% Generated from sealed predictions; all positive windows retained.
\begin{table}[t]\centering
\caption{Retrospective RAW family-and-source correctness on every physical-positive window, using the archived source metadata and current event-code support. Multiple windows from a sustained annotation are not independent actions.}
\label{tab:raw-source-hybrid}
\small\setlength{\tabcolsep}{5pt}
\begin{tabular}{lrr}\hline
Family & Submitted ZIP & Source hybrid\\\hline
Fault & 1/1 & 1/1 \\
Line & 0/1 & 1/1 \\
Generation & 0/11 & 0/11 \\
Load & 0/11 & 0/11 \\
Total & 1/24 & 2/24 \\
\hline\end{tabular}
\end{table}

Joint correctness changes from \RawZIPJoint{}/\RawZIPTotal{} to
\RawHybridJoint{}/\RawHybridTotal{}: \RawHybridGains{} paired gain and
\RawHybridLosses{} losses, with \RawHybridFallbacks{}/\SSRawWindows{}
total windows retaining the ZIP source. The gain corrects the line location; load and generation
remain unresolved. A single improved window in one previously used record
does not establish reliable all-family transfer.

The implemented interface combines these physical and integrity outputs in
one chronological stream. Exact projection checks recover all
\SystemWindows{} component windows and the availability state across
\SystemReports{} reports. This verifies composition without adding a
performance gain; completed-window evidence cannot reopen an already
restored missing-PMU episode.

The calendar
recovered for the archived hackathon table does not match the actual RAW
annotations; its score is not a reproduced comparator. Relabeling windows
or relaxing one admission gate would not validate the physical/instrument
model. The archive preserves these failures and the current-calendar
evaluation. Superiority over the
historical all-family system remains unestablished.
